\capitulo{4}{Técnicas y herramientas}

Este capítulo tiene como finalidad describir, de manera ordenada y argumentada, las técnicas metodológicas y las herramientas empleadas a lo largo del desarrollo del proyecto.

\section{Técnicas metodológicas}

A lo largo del desarrollo se han aplicado técnicas orientadas a organizar el trabajo, mantener consistencia en el código y facilitar la evolución incremental del prototipo. En concreto, se ha adoptado una metodología de planificación inspirada en \textit{Scrum}, se han seguido convenciones de estilo para mejorar la legibilidad y mantenibilidad del código, y se ha aplicado un enfoque sistemático de construcción de interfaz mediante un sistema de utilidades, con el fin de asegurar coherencia visual sin introducir deuda técnica innecesaria.

\subsection{Planificación iterativa basada en Scrum}

La planificación se ha estructurado en iteraciones, denominadas \textit{sprints}, con el objetivo de descomponer el alcance del proyecto en tareas abordables y reducir el riesgo de bloqueo en fases finales. Para cada iteración se han definido objetivos concretos, priorizados en función del valor aportado al sistema y de sus dependencias técnicas, manteniendo una visión continua del trabajo pendiente mediante un \textit{backlog}. Este enfoque ha permitido combinar desarrollo funcional y redacción de memoria de forma coordinada, incorporando revisiones periódicas del avance y ajustes razonables cuando se detectaban desajustes entre estimación y esfuerzo real.

\subsection{Convenciones de estilo y guía PEP8}

Para asegurar consistencia y facilitar la revisión del código, se han seguido convenciones de estilo alineadas con PEP8 en los módulos Python del proyecto. Esta decisión contribuye a que el código sea más homogéneo, reduce fricción al introducir nuevas funcionalidades y mejora la mantenibilidad del repositorio, especialmente en zonas críticas como la gestión de rutas, la persistencia de resultados y la integración del \textit{pipeline} de inferencia. Además, favorece la detección temprana de errores comunes y simplifica la colaboración y supervisión del proyecto por parte de terceros.

\subsection{Diseño de interfaz con Tailwind CSS y DaisyUI}

El diseño de la interfaz se ha abordado mediante un enfoque basado en utilidades con \textit{Tailwind CSS}, complementado con \textit{DaisyUI} como capa de componentes. Este planteamiento permite construir vistas consistentes a partir de clases de bajo nivel, manteniendo un control fino sobre espaciados, tipografía y jerarquía visual, y evitando el coste de sobrescrituras frecuentes asociado a estilos rígidos. En la práctica, el uso de un sistema de utilidades favorece la escalabilidad del \textit{frontend}, ya que la coherencia visual emerge de patrones reutilizables y de una configuración centralizada, lo que resulta especialmente útil en una aplicación con pantallas heterogéneas, como procesamiento de imágenes, visor cartográfico y paneles de administración.

\subsection{Internacionalización e independencia del idioma}

El sistema se ha diseñado con capacidad de internacionalización, desacoplando los textos visibles del idioma concreto mediante el uso de catálogos de traducción. Esta técnica permite que la aplicación sea neutral desde el punto de vista lingüístico, evitando duplicidad de código y facilitando añadir nuevos idiomas sin modificar la lógica principal. En el caso del contenido dinámico generado en cliente, se sigue una estrategia consistente para que los textos sigan siendo localizables y no queden incrustados de manera rígida en el \textit{frontend}.

\subsection{Gestión de dependencias, reproducibilidad y calidad}

Para garantizar reproducibilidad, el entorno de ejecución se define mediante un fichero de dependencias versionadas y un entorno aislado de Python. Esta decisión reduce problemas derivados de incompatibilidades entre versiones, facilita replicar el entorno de desarrollo y estabiliza el despliegue. En paralelo, se siguen criterios de estilo y organización del código, manteniendo convenciones consistentes y promoviendo cambios incrementales. Finalmente, el proyecto incorpora una estrategia de pruebas automatizadas que permite detectar regresiones en funcionalidades críticas, especialmente en aquellas que integran inferencia, persistencia y comunicación cliente--servidor.

\section{Herramientas y tecnologías empleadas}

En esta sección se describen las herramientas principales utilizadas para implementar el sistema, agrupadas según su función dentro del \textit{stack}. La finalidad no es enumerar exhaustivamente cada dependencia, sino destacar aquellas que han sido determinantes para construir la aplicación, mantener un desarrollo sostenible y asegurar reproducibilidad.

\subsection{Entorno de ejecución y dependencias}

\begin{itemize}
	\item \textbf{Python:} lenguaje base del proyecto.
	\item \textbf{\texttt{pip} y \texttt{requirements.txt}:} instalación y fijación de dependencias para reproducir el entorno.
	\item \textbf{Entorno virtual y \textit{conda}:} mecanismos de aislamiento del entorno, empleados para evitar colisiones de dependencias y estabilizar la ejecución.
\end{itemize}

\subsection{Backend y capa web}

\begin{itemize}
	\item \textbf{Flask:} \textit{microframework} utilizado para implementar el servidor, rutas y lógica de aplicación.
	\item \textbf{Jinja2:} motor de plantillas para renderizado del HTML en servidor.
	\item \textbf{Werkzeug:} utilidades WSGI y soporte de infraestructura del ecosistema Flask.
	\item \textbf{Flask-Login:} gestión de autenticación, sesión y permisos por rol.
	\item \textbf{Flask-WTF y WTForms:} formularios y validación robusta de entradas.
	\item \textbf{Flask-Babel y Babel:} internacionalización y localización de la interfaz.
	\item \textbf{Flask-Admin:} soporte para vistas y operaciones de administración.
\end{itemize}

\subsection{Persistencia y acceso a datos}

\begin{itemize}
	\item \textbf{SQLite:} base de datos embebida utilizada para persistencia durante el desarrollo.
	\item \textbf{SQLAlchemy y Flask-SQLAlchemy:} capa ORM para desacoplar el acceso a datos y facilitar mantenibilidad.
\end{itemize}

\subsection{Inferencia, modelos y procesamiento de imagen}

\begin{itemize}
	\item \textbf{PyTorch:} marco principal para ejecutar el modelo de segmentación en inferencia.
	\item \textbf{\texttt{pickle}:} mecanismo de deserialización para modelos heredados.
	\item \textbf{\texttt{segmentation-models-pytorch}:} soporte para arquitecturas y componentes de segmentación.
	\item \textbf{\texttt{timm}:} catálogo de \textit{backbones} y modelos preentrenados empleados como codificadores.
	\item \textbf{\texttt{safetensors}:} formato seguro y eficiente para el almacenamiento y carga de pesos de modelos cuando procede.
	\item \textbf{Jupyter Notebook:} entorno utilizado para explorar el \textit{pipeline} original, replicar inferencias y validar resultados antes de integrarlos en la aplicación.
	\item \textbf{NumPy y Pillow:} utilidades para manipulación de arrays y lectura o conversión de imágenes.
\end{itemize}

\subsection{Frontend, estilos y experiencia de usuario}

\begin{itemize}
	\item \textbf{HTML5, CSS y JavaScript:} tecnologías base del \textit{frontend}.
	\item \textbf{Tailwind CSS:} \textit{framework} de utilidades empleado para el diseño y la composición visual.
	\item \textbf{DaisyUI:} extensión basada en componentes sobre Tailwind, utilizada para acelerar la construcción de elementos recurrentes manteniendo coherencia.
\end{itemize}

\subsection{Cartografía y datos geoespaciales}

\begin{itemize}
	\item \textbf{Leaflet:} biblioteca de mapas en cliente utilizada para el visor interactivo.
	\item \textbf{OpenStreetMap:} mapa base de referencia para navegación e interacción.
	\item \textbf{PNOA mediante servicios WMS:} fuente de ortofotografía de alta resolución utilizada para visualización y obtención de recortes.
\end{itemize}

\subsection{Pruebas, diagramación y documentación}

\begin{itemize}
	\item \textbf{pytest:} \textit{framework} de pruebas automatizadas para reducir regresiones.
	\item \textbf{Git y GitHub:} control de versiones y gestión del repositorio.
	\item \textbf{Visual Studio Code:} editor principal de desarrollo.
	\item \textbf{Diagrams.net (Draw.io):} herramienta empleada para la elaboración de diagramas de arquitectura y flujo.
	\item \textbf{\LaTeX{}, \TeX{}Studio y MiK\TeX{} Console:} herramientas utilizadas para redactar, compilar y gestionar la memoria y sus dependencias tipográficas.
\end{itemize}

\subsection{Gestión del proyecto}
\begin{itemize}
	\item \textbf{Zube.io:} herramienta utilizada para gestionar el \textit{backlog} y organizar las tareas por \textit{sprints}, facilitando el seguimiento del avance y la priorización del trabajo.
\end{itemize}

\subsection{Comparativa entre Bootstrap y Tailwind CSS}

En el diseño e implementación de la interfaz web del proyecto resulta necesario seleccionar un \textit{framework} de estilos que aporte productividad sin comprometer la mantenibilidad ni la calidad del código. Entre las alternativas más consolidadas se encuentran Bootstrap, orientado a componentes predefinidos, y Tailwind CSS, basado en un paradigma de utilidades. A continuación se presenta una comparativa crítica entre ambas opciones, justificando la elección final de Tailwind CSS para el desarrollo de esta página.

Bootstrap propone un conjunto amplio de componentes listos para usar, tales como barras de navegación, tarjetas o sistemas de rejilla, acompañados de un diseño visual relativamente homogéneo. Esta aproximación permite, en un primer momento, acelerar el desarrollo, ya que se pueden componer interfaces funcionales mediante la combinación de bloques estándar. Además, Bootstrap ofrece una documentación madura y una comunidad extensa, lo que facilita la resolución de dudas y la búsqueda de ejemplos; sin embargo, esta misma estandarización introduce ciertas limitaciones. Por un lado, la personalización profunda del diseño suele requerir sobrescribir reglas de estilo, lo que incrementa la complejidad de las hojas CSS y genera deuda técnica con facilidad. Por otro lado, las interfaces tienden a parecerse entre sí, de forma que la identidad visual del proyecto puede quedar diluida si no se realiza un esfuerzo considerable de adaptación.

En contraste, Tailwind CSS se basa en clases utilitarias de bajo nivel que describen propiedades concretas, tales como márgenes, tipografía o colores, aplicadas directamente sobre el marcado HTML. Esta filosofía permite construir componentes a medida a partir de combinaciones de utilidades, sin depender de un conjunto rígido de componentes predefinidos. Como consecuencia, el diseño resultante se ajusta mejor a los requisitos específicos del proyecto y favorece la creación de un sistema de diseño propio, más alineado con la identidad visual deseada. Además, Tailwind CSS integra mecanismos de purga de clases no utilizadas, lo que reduce de manera significativa el tamaño del fichero CSS final y mejora el rendimiento en producción.

No obstante, Tailwind CSS presenta también ciertos inconvenientes. Al inicio, la curva de aprendizaje puede parecer más pronunciada, ya que obliga a pensar en términos de propiedades de bajo nivel y no tanto en componentes cerrados. Asimismo, la presencia de múltiples clases en cada elemento puede dar lugar a un marcado aparentemente más denso, pero, cuando se interioriza el conjunto de utilidades y se adoptan buenas prácticas, como la extracción de componentes reutilizables y el uso de configuración centralizada, el resultado es un código más predecible y coherente, con menos necesidad de sobrescrituras complejas.

Desde una perspectiva de la ingeniería del software, la elección entre ambos enfoques afecta de forma directa a la mantenibilidad y escalabilidad del proyecto. Bootstrap ofrece rapidez inicial, pero puede derivar en hojas de estilo difíciles de gestionar cuando se requiere un alto grado de personalización. En cambio, Tailwind CSS promueve una relación más estrecha entre el diseño y el código, ya que el sistema de diseño se expresa mediante utilidades consistentes que se reutilizan en toda la aplicación. Esto reduce la duplicidad de estilos, minimiza el riesgo de efectos colaterales al modificar clases existentes y facilita la evolución del producto a largo plazo.

Finalmente, he optado por Tailwind CSS debido a su capacidad para generar interfaces más ligeras, a su enfoque modular y a su mejor integración con patrones modernos de desarrollo basado en componentes. Gracias a su sistema de utilidades, ha sido posible definir un diseño adaptado a las necesidades concretas del proyecto, mantener un control fino sobre el \textit{bundle} de estilos y reducir la deuda técnica asociada a la personalización del \textit{frontend}. En definitiva, Tailwind CSS proporciona un equilibrio más adecuado entre flexibilidad, rendimiento y mantenibilidad, lo que lo convierte en la opción más idónea para este desarrollo.

\begin{table}[h]
	\centering
	\renewcommand{\arraystretch}{1.3}
	\begin{tabular}{p{0.18\textwidth}p{0.36\textwidth}p{0.36\textwidth}}
		\toprule
		\textbf{Aspecto} & \textbf{Bootstrap} & \textbf{Tailwind CSS} \\
		\midrule
		Enfoque & Conjunto de componentes predefinidos y estilos de alto nivel &
		Clases utilitarias de bajo nivel que permiten construir componentes a medida \\
		Curva de aprendizaje & Accesible al inicio gracias a los componentes estándar y a los ejemplos completos &
		Inicialmente más exigente, requiere pensar en términos de propiedades y sistemas de diseño \\
		Productividad inicial & Elevada en prototipos y desarrollos con requisitos visuales genéricos &
		Crece con el tiempo, especialmente cuando se definen patrones reutilizables y componentes propios \\
		Personalización del diseño & Personalización profunda costosa, necesidad frecuente de sobrescribir estilos y romper el diseño base &
		Alta capacidad de personalización mediante combinaciones de utilidades y configuración centralizada \\
		Mantenibilidad & Riesgo de hojas CSS voluminosas y difíciles de refactorizar en proyectos muy adaptados &
		Estilos más controlados, menor necesidad de sobrescrituras y reducción de la deuda técnica \\
		Rendimiento y tamaño del CSS & Tamaño mayor en el fichero final cuando se utilizan muchos componentes sin optimización específica &
		Tamaño significativamente reducido gracias a la purga de utilidades no utilizadas en el producto final \\
		Identidad visual & Interfaces con aspecto reconocible y similar a otros proyectos basados en el mismo \textit{framework} &
		Mayor capacidad para definir una identidad visual propia y diferenciada \\
		Comunidad y ecosistema & Comunidad muy amplia, gran cantidad de plantillas, ejemplos y recursos formativos &
		Comunidad en fuerte crecimiento, ecosistema consolidado en proyectos modernos basados en componentes \\
		\bottomrule
	\end{tabular}
	\caption{Comparativa de ventajas e inconvenientes entre Bootstrap y Tailwind CSS.}
	\label{tab:bootstrap_tailwind}
\end{table}

\subsection{Comparativa de motores y fuentes de datos geoespaciales}

La elección del motor geoespacial y de la fuente de datos constituye una decisión crítica en el diseño de aplicaciones web orientadas al análisis terrenal y a la detección de elementos lineales, ya que condiciona de forma directa tanto el nivel de detalle alcanzable como la escalabilidad, la reproducibilidad y la viabilidad legal del sistema. En este contexto, se han analizado cinco alternativas representativas del estado del arte actual, valorando sus capacidades técnicas, sus limitaciones y su adecuación a los objetivos del proyecto.

En primer lugar, Google Earth Engine se presenta como una plataforma de computación geoespacial en la nube que destaca por su amplio catálogo de datos científicos y por su capacidad de procesado a gran escala. Permite trabajar de forma eficiente con series temporales extensas y con colecciones como Sentinel-2 o Landsat, facilitando la generación de mosaicos, máscaras de nubes o índices espectrales. No obstante, su principal limitación en este proyecto radica en la resolución espacial de dichas fuentes, que suele resultar insuficiente para la detección de trazas finas, como senderos o caminos estrechos. Además, conviene subrayar que la imaginería de alta resolución visible en Google Earth no forma parte de los datasets públicos accesibles para análisis y exportación, lo que restringe su uso como fuente principal de detalle.

Como alternativa orientada a la integración en aplicaciones web, Sentinel Hub y el Copernicus Data Space ofrecen interfaces de programación que permiten solicitar imágenes procesadas bajo demanda. Esta aproximación resulta especialmente cómoda en arquitecturas cliente-servidor, ya que simplifica el flujo de petición y respuesta mediante recortes espaciales y temporales. Sin embargo, al igual que en el caso anterior, la resolución de las imágenes Sentinel constituye un factor limitante cuando se requiere un nivel de detalle submétrico, lo que reduce su idoneidad para el objetivo concreto del proyecto.

Una tercera opción, de especial relevancia en el ámbito nacional, es el uso de ortofotos aéreas de alta resolución, destacando en España el Plan Nacional de Ortofotografía Aérea, PNOA. Este conjunto de datos proporciona imágenes con resoluciones del orden de decímetros, muy superiores a las de los sensores satelitales de libre acceso. Además, su disponibilidad a través de servicios WMS y WMTS, junto con opciones de descarga directa, lo convierte en una solución sólida tanto para la visualización como para el análisis reproducible. Desde una perspectiva académica, el uso de fuentes oficiales y abiertas refuerza la validez metodológica del trabajo y evita problemas derivados de licencias restrictivas.

En contraposición, el uso directo de imaginería procedente de Google Maps Imagery o Google Earth como fuente de datos para análisis automático presenta importantes inconvenientes. A pesar de su elevada calidad visual, estas plataformas imponen restricciones estrictas en cuanto a la extracción, almacenamiento y procesamiento de imágenes, lo que dificulta su empleo como motor principal en un proyecto de carácter académico y analítico. Por este motivo, su utilización queda prácticamente relegada a contextos de visualización, no siendo recomendable como base para flujos de detección y análisis a escala.

Finalmente, en lo relativo a la capa de visualización, se ha considerado el uso de Folium frente a la integración directa de bibliotecas JavaScript como Leaflet. Aunque Folium resulta adecuado para prototipos rápidos generados desde Python, su naturaleza estática y su dependencia del backend lo hacen menos apropiado para aplicaciones web interactivas complejas. En consecuencia, una solución basada en Leaflet u OpenLayers en el frontend, combinada con un backend que sirva datos vectoriales en formato JSON, ofrece mayor flexibilidad y control.

A la luz de este análisis comparativo, la solución adoptada en el proyecto consiste en utilizar ortofotografía aérea de alta resolución, concretamente PNOA, como motor principal de datos geoespaciales. Esta elección garantiza el nivel de detalle necesario para la detección de trazas finas, asegura la reproducibilidad del trabajo y se ajusta a criterios de legalidad y sostenibilidad académica. De forma complementaria, se contempla el uso de Google Earth Engine como herramienta de preprocesado o filtrado a gran escala, aprovechando su potencia para análisis masivos sin depender de él como fuente final de detalle. Este enfoque híbrido proporciona un equilibrio óptimo entre precisión espacial, escalabilidad y rigor metodológico.

\begin{table}[h]
	\centering
	\renewcommand{\arraystretch}{1.3}
	\begin{tabular}{p{3 cm} p{4 cm} p{5 cm}}
		\hline
		\textbf{Alternativa} & \textbf{Ventajas principales} & \textbf{Limitaciones} \\
		\hline
		Google Earth Engine & Gran catálogo de datos científicos y alto rendimiento en análisis masivo & Resolución insuficiente para trazas finas y ausencia de imaginería Google Earth exportable \\
		\hline
		Sentinel Hub / Copernicus & Integración sencilla mediante API y procesado bajo demanda & Resolución limitada a sensores Sentinel y dependencia de cuotas \\
		\hline
		Ortofoto aérea (PNOA) & Alta resolución submétrica, fuentes oficiales y buena reproducibilidad & Cobertura limitada al ámbito nacional \\
		\hline
		Google Maps / Earth imagery & Excelente calidad visual para visualización & Restricciones de licencia y uso no adecuado para análisis automático \\
		\hline
		Folium / Leaflet & Facilidad en prototipos desde Python & Menor flexibilidad para aplicaciones web interactivas avanzadas \\
		\hline
	\end{tabular}
	\caption{Comparativa de motores y fuentes de datos geoespaciales}
	\label{tab:comparativa_engines}
\end{table}

\section{Comparativa de extensiones para la gestión de bases de datos en Visual Studio Code}

Para el desarrollo de este proyecto la elección de la herramienta utilizada para la gestión y exploración de la base de datos resulta especialmente relevante. En este Trabajo Fin de Grado, dicha elección no se limita únicamente al sistema gestor de bases de datos, sino también a la extensión empleada dentro del entorno de desarrollo integrado. Una herramienta adecuada debe facilitar la inspección de esquemas, la ejecución de consultas y la depuración de datos, sin introducir una complejidad innecesaria ni dependencias excesivas.

Entre las alternativas más consolidadas del estado del arte actual destaca, en primer lugar, SQLTools, una extensión de carácter modular que actúa como núcleo común al que se añaden controladores específicos según el motor de base de datos utilizado. Este enfoque permite trabajar de forma homogénea con diferentes sistemas, como SQLite, PostgreSQL o MySQL, sin modificar la herramienta principal. Desde un punto de vista arquitectónico, esta solución resulta especialmente atractiva en proyectos que pueden evolucionar con el tiempo, ya que evita la dependencia temprana de un único motor. No obstante, esta flexibilidad conlleva una ligera sobrecarga conceptual, puesto que el usuario debe gestionar explícitamente la instalación y configuración de los controladores necesarios.

Como alternativa especializada, la extensión SQLite, desarrollada por alexcvzz, está orientada exclusivamente a la gestión de este tipo de bases de datos, caracterizadas por funcionar como un único archivo local y no requerir de un servidor dedicado. Su propuesta prioriza la simplicidad y la inmediatez, ofreciendo un explorador visual de tablas, ejecución directa de sentencias SQL y opciones de exportación de resultados, todo ello con una configuración mínima. Esta aproximación resulta especialmente adecuada en fases iniciales del desarrollo, donde la base de datos se materializa como un fichero local y se requiere una herramienta ligera que no interfiera en el flujo de trabajo. La principal limitación de esta extensión radica en su carácter específico, ya que una eventual migración a otro motor implicaría adoptar una herramienta diferente.

Otra opción relevante es Database Client de Weijan Chen, una extensión que aspira a funcionar como un cliente de bases de datos completo dentro de Visual Studio Code, soportando múltiples motores e incluso conexiones avanzadas mediante túneles SSH. Aunque su nivel de funcionalidades es elevado, existen consideraciones importantes desde el punto de vista académico, como su transición hacia un modelo parcialmente cerrado y la presencia de mecanismos de telemetría. En el contexto de un Trabajo Fin de Grado, donde se valora la transparencia de las herramientas y la reproducibilidad del entorno, estos aspectos reducen su idoneidad como opción principal.

Finalmente, la extensión PostgreSQL for VS Code de ms-ossdata ofrece un conjunto de herramientas específicas para trabajar con PostgreSQL, tanto en entornos locales como en despliegues en la nube. Se trata de una solución robusta cuando el proyecto tiene claramente definido este motor como destino final. Sin embargo, su especialización la hace poco adecuada en fases tempranas del desarrollo, especialmente cuando se trabaja con SQLite como base de datos embebida, ya que obliga a combinar múltiples extensiones para cubrir distintas necesidades.

Tras analizar las distintas alternativas, se ha optado por emplear SQLite como sistema gestor de bases de datos durante el desarrollo del proyecto, apoyándose en una extensión ligera y especializada para su gestión dentro de Visual Studio Code. Esta decisión responde a varios criterios fundamentales. En primer lugar, SQLite permite trabajar con una base de datos autocontenida, sin necesidad de desplegar servicios adicionales ni configurar infraestructuras externas, lo que simplifica el entorno de desarrollo y reduce la fricción inicial. En segundo lugar, su integración natural con aplicaciones web ligeras y con frameworks como Flask lo convierte en una solución especialmente adecuada para proyectos académicos de alcance controlado.

Desde la perspectiva del entorno de desarrollo, el uso de una extensión específica para SQLite ofrece una experiencia directa, coherente y alineada con los objetivos del proyecto, facilitando la inspección de datos y la validación del modelo relacional sin introducir complejidades innecesarias. Además, esta elección no compromete la evolución futura del sistema, ya que el diseño de la capa de persistencia permite, llegado el caso, una migración hacia motores más robustos como PostgreSQL, apoyándose entonces en herramientas más generalistas.

En consecuencia, la combinación de SQLite como motor de base de datos y una extensión dedicada en Visual Studio Code representa la opción con mejor equilibrio entre simplicidad y eficiencia, garantizando un desarrollo controlado y sostenible.

\begin{table}[h]
	\centering
	\renewcommand{\arraystretch}{1.3}
	\begin{tabular}{p{0.28\textwidth} p{0.3\textwidth} p{0.32\textwidth}}
		\hline
		\textbf{Extensión} & \textbf{Ventajas principales} & \textbf{Limitaciones} \\
		\hline
		SQLTools & Modular, soporta múltiples motores y facilita la evolución del proyecto & Requiere gestión de controladores y mayor configuración inicial \\
		\hline
		SQLite (alexcvzz) & Simplicidad, integración directa con ficheros SQLite y bajo coste cognitivo & Limitada a SQLite, no válida para otros motores \\
		\hline
		Database Client & Muy completa y con soporte para múltiples tecnologías & Modelo parcialmente cerrado y menor transparencia \\
		\hline
		PostgreSQL for VS Code & Herramientas específicas y maduras para PostgreSQL & Inadecuada para trabajar con SQLite \\
		\hline
	\end{tabular}
	\caption{Comparativa de extensiones para la gestión de bases de datos en Visual Studio Code}
	\label{tab:comparativa_extensiones_bd}
\end{table}


